\exnumber{01}
\extitle{Projet 4 : À table !}
\exfiles{ex01/archive zip, nomme Name\_FirstName\_DeliverableName\_Date\_VersionNumber}
\exauthorize{Those named in the exercise}
\makeheaderfiles

Show a tested high-fidelity prototype of each of the 4 journeys.

\textbf{Estimated time:} 5 weeks

\textbf{Team size:} max 3 people

\textbf{Consignes :}

\begin{itemize}
\item Testez vos prototypes avec au moins 5 personnes.
\item Make your presentations as if you were showing your work to DeliverNow team. Make sure there are at least 5 people in the audience.
\item Préparez un formulaire de retour pour le public. Il contient au moins les questions listées plus bas.
\item Timing: 15 minutes for the presentation and 5 minutes for Q\&A.
\item Vous pouvez vous servir des audits de présentation pour structurer vos présentations.
\item Déposez les documents sur Github, dans un zip nommé « ProjectTitle\_Name\_FirstName », et chaque livrable ainsi : « Name\_FirstName\_DeliverableName\_Date\_VersionNumber ». Exemple : « Doe\_John\_MidFiPrototype\_05242024\_V1 ».
\item Ajoutez un titre dans tous les documents écrits.
\item Soignez l'aspect global des livrables. Ils doivent être clairs, simples et lisibles. Vous pouvez regarder canva.com, sans surcharger le dessin de détails.
\item Amusez-vous !
\end{itemize}

\textbf{Outils :}

\begin{itemize}
\item Google Forms \url{https://docs.google.com/forms/u/0/}
\item Airtable \url{https://www.airtable.com/}
\item Typeform \url{https://www.typeform.com/}
\end{itemize}

\textbf{Pistes :}

\begin{itemize}
\item Effective presentation skills \url{https://academic.oup.com/femsle/article/364/24/fnx235/4587905}
\end{itemize}

Livrables intermédiaires :

\textbf{Citation :}

« Your mission in any presentation is to inform, educate, and inspire. » Robert Ballard

\medskip

\textbf{Phase 1 : audit de la présentation intermédiaire}

Vérifiez :

\begin{itemize}
\item Au moins 5 personnes assistent à la présentation
\item La présentation porte sur tout le processus UX, y compris :
\begin{itemize}
\item Le brief initial
\item Un planning, avec les phases UX, les outils choisis, le calendrier
\item User research findings from qualitative and quantitative research (persona, figures, etc.) for the 4 users
\item 4 problem statements to adapt the initial brief based on the research findings
\item An introduction of each concept that came out of the ideation
\item The animated mid-fidelity prototype (a demo-video) with the 4 user journeys
\item The main feedback from usability testings (5x4)
\item Les prochaines étapes du projet
\end{itemize}
\item La présentation ne dépasse pas 15 minutes
\item Les formulaires de retour ont été distribués au début de la présentation
\item Vous vous êtes amusés !
\end{itemize}

\medskip

\textbf{Phase 2 : audit de la présentation finale}

Vérifiez :

\begin{itemize}
\item Au moins 5 personnes assistent à la présentation
\item La présentation porte sur le processus UI, y compris :
\begin{itemize}
\item Un bref rappel de la présentation précédente
\item Un planning, avec les phases UI, les outils choisis, le calendrier
\item Une planche d'ambiance
\item Un styleguide (couleurs, marque, typographie, quelques éléments de la bibliothèque)
\item The animated high-fidelity prototype (a demo-video) with the 4 user journeys
\item Les retours principaux des tests d'utilisabilité et de désirabilité
\item Les prochaines étapes du projet
\end{itemize}
\item La présentation ne dépasse pas 15 minutes
\item Les formulaires de retour ont été distribués au début de la présentation
\item Vous vous êtes amusés !
\end{itemize}

\textbf{Questions à poser dans les formulaires de retour}

\begin{itemize}
\item Que retenez-vous de cette présentation ?
\item De 0 à 10, le discours était-il clair ?
\item Pourquoi ?
\item De 0 à 10, la présentation était-elle engageante ?
\item Pourquoi ?
\item De 0 à 10, la présentation était-elle cohérente ?
\item Pourquoi ?
\item De 0 à 10, les visuels portaient-ils ?
\item Pourquoi ?
\item De 0 à 10, le problème était-il clair ?
\item Pourquoi ?
\item De 0 à 10, le prototype répond-il au problème ?
\item Pourquoi ?
\item Quelle recommandation feriez-vous pour améliorer cette présentation ?
\item Pourquoi ?
\end{itemize}

